% ============================================================
%  TOGT --- Volume Two: Outline
%  Renormalization, Topological Invariants, and the Crystalline Spine
%  Author: Pablo Grossi
%  Date:   March 2026
% ============================================================
\documentclass[11pt]{article}

\usepackage[T1]{fontenc}
\usepackage[utf8]{inputenc}
\usepackage{geometry}
\geometry{margin=1in}
\usepackage{amsmath,amssymb}
\usepackage{xcolor}
\usepackage[colorlinks=true,linkcolor=blue,citecolor=purple,urlcolor=teal]{hyperref}
\usepackage{enumitem}

\emergencystretch=1em

% Status markers
\newcommand{\provable}{\textcolor{blue!70!black}{\textbf{[Provable target]}}}
\newcommand{\conjectural}{\textcolor{purple}{\textbf{[Conjectural]}}}
\newcommand{\depends}[1]{\smallskip\noindent\textit{Depends on (V1):} #1.\smallskip}

\newcommand{\g}{\boldsymbol{g}}
\newcommand{\Comp}{\mathcal{C}\!omp}
\newcommand{\Konst}{\mathcal{K}\!onst}
\newcommand{\Fold}{\mathcal{F}\!old}
\newcommand{\Unfold}{\mathcal{U}\!nfold}

\title{Topographical Orthogonal Generative Theory (TOGT)\\[4pt]
\large Volume Two: Renormalization, Topological Invariants,\\
and the Crystalline Spine\\[6pt]
\normalsize \emph{Structural Outline}}
\author{Pablo Grossi}
\date{March 2026}

\begin{document}
\maketitle

\begin{abstract}
This document is a structural outline for Volume Two of TOGT. It organises
the four open research directions identified in Volume One into a coherent
sequence of sections, specifying dependencies, provable targets, and
conjectural directions. No proofs or mathematical claims are made here;
the purpose is architectural scaffolding only.

Two typographic markers are used throughout:
\provable{} marks goals for which a proof strategy is in view and the
mathematical prerequisites are in place;
\conjectural{} marks directions that require new definitions, external
formalisation, or unresolved foundational questions before a proof can be
attempted.
\end{abstract}

\tableofcontents
\bigskip

% ============================================================
\part{Function-Space Foundations for the Renormalization Operator}
% ============================================================

%---------------------------------------------------------------
\section{Unimodal Map Spaces and Their Norms}
%---------------------------------------------------------------

\subsection*{Description}
This section establishes the precise function space in which the
renormalization operator $\mathcal{R}$ will be defined. The two natural
candidates --- analytic unimodal maps on a complex neighbourhood of
$[-1,1]$, and $C^3$ unimodal maps with negative Schwarzian derivative ---
are compared, and the relevant Banach or Fr\'{e}chet norms are recorded.
The section closes by fixing one setting as the working hypothesis for the
rest of Volume Two, with a justification based on the regularity of the
Euler-discretized gradient flow $f_{\lambda,h}$.

\subsection*{Status}
\provable{} The definitions and norm comparisons are standard material from
the Collet--Eckmann and Sullivan--McMullen literature and require no new
mathematics; they are imported and specialised to the TOGT setting.

\depends{Proposition~3.2 (V1, \S3.3): the Euler discretization
$f_{\lambda,h}$ is a $C^\infty$ unimodal map with cubic critical point,
hence lies in both candidate spaces for $h$ small enough}

%---------------------------------------------------------------
\section{The Renormalization Operator as a Map on Function Space}
%---------------------------------------------------------------

\subsection*{Description}
Having fixed the function space, this section defines $\mathcal{R}$ as a
rigorous operator
\[
  \mathcal{R} : \mathcal{U} \to \mathcal{U},
  \qquad
  \mathcal{R}(f) = \alpha^{-1} \circ f^2 \circ \alpha,
\]
where $\mathcal{U}$ is the chosen space of unimodal maps and $\alpha$ is
the Feigenbaum rescaling. Domain invariance --- the question of whether
$\mathcal{R}(f)$ remains unimodal when $f$ is --- is verified. Basic
continuity and compactness properties are established.

\subsection*{Status}
\provable{} Domain invariance and continuity follow from the classical theory
(Collet--Eckmann~\cite{collet1980}); the contribution here is specialising
those results to the TOGT operator family and verifying that
$f_{\lambda,h}$ lies in the domain.

\depends{Section~1 of this volume (function-space choice);
Proposition~3.2 (V1): regularity of $f_{\lambda,h}$}

%---------------------------------------------------------------
\section{Relation Between the Gradient Flow and Its Discretization}
%---------------------------------------------------------------

\subsection*{Description}
The gradient flow $\g_\lambda$ itself is not a unimodal map and does not
lie in $\mathcal{U}$; Volume One states this explicitly (\S3.3). This
section makes the relationship precise: it characterises the map
$h \mapsto f_{\lambda,h}$ and studies whether the renormalization orbit of
$f_{\lambda,h}$ converges as $h \to 0$ to a well-defined limit that can be
attributed to the continuous flow. This is the key step needed to connect
$\mathcal{R}$ to the TOGT pipeline proper.

\subsection*{Status}
\conjectural{} The $h \to 0$ limit of the renormalization orbit is not
covered by existing results. A rigorous treatment requires either a
continuous-time renormalization theory (non-standard) or a controlled
discretisation argument. This is the primary new mathematical challenge of
Part~I.

\depends{Definition of $\g_\lambda$ (V1, \S3.1);
Proposition~3.2 (V1): domain membership of $f_{\lambda,h}$;
Section~2 of this volume: $\mathcal{R}$ well-defined on $\mathcal{U}$}

% ============================================================
\part{Dynamics of the Renormalization Operator}
% ============================================================

%---------------------------------------------------------------
\section{The Feigenbaum Fixed Point: Existence and Uniqueness}
%---------------------------------------------------------------

\subsection*{Description}
This section imports the Feigenbaum--Coullet--Tresser fixed-point theorem
(as proved analytically by Lanford, and later by Lyubich) into the TOGT
setting. The fixed point $\g_\infty \in \mathcal{U}$ of $\mathcal{R}$ is
identified, and its basic properties --- symmetry, analyticity, and the
values of the universal constants $\alpha \approx 2.5029$ and
$\delta \approx 4.6692$ --- are recorded. The section does not re-prove
the fixed-point theorem; it establishes that the TOGT operator
$f_{\lambda,h}$ lies in the basin of attraction.

\subsection*{Status}
\provable{} The fixed-point theorem is established in the literature;
the work here is verifying the basin condition for $f_{\lambda,h}$, which
follows from Proposition~3.2 (V1) and the universality class membership.

\depends{Proposition~3.2 (V1): $f_{\lambda,h}$ lies in the
period-doubling universality class;
Section~2 of this volume: $\mathcal{R}$ well-defined}

%---------------------------------------------------------------
\section{Hyperbolicity of the Renormalization Horseshoe}
%---------------------------------------------------------------

\subsection*{Description}
The renormalization operator $\mathcal{R}$ has a hyperbolic fixed point
$\g_\infty$ with a one-dimensional unstable manifold (corresponding to
the quadratic family parameter) and an infinite-codimensional stable
manifold (the space of infinitely renormalizable maps of the same
combinatorial type). This section states the Avila--Lyubich hyperbolicity
theorem precisely and identifies its implications for the TOGT pipeline:
in particular, which parameters $(\lambda, h)$ yield infinitely
renormalizable maps.

\subsection*{Status}
\provable{} Hyperbolicity is established by Avila--Lyubich~\cite{avila2011};
the contribution is computing or bounding the set of $(\lambda, h)$
for which $f_{\lambda,h}$ is infinitely renormalizable.

\depends{Section~4 of this volume: fixed-point identification;
the contraction and pitchfork results (V1, \S\S5--6): they constrain
the parameter range of interest}

%---------------------------------------------------------------
\section{Stable Manifolds and the Basin of the Universal Fixed Point}
%---------------------------------------------------------------

\subsection*{Description}
The stable manifold of $\g_\infty$ under $\mathcal{R}$ consists of all
infinitely renormalizable maps with the same combinatorial type (period
doubling). This section characterises this manifold within $\mathcal{U}$,
identifies which maps $f_{\lambda,h}$ lie on it, and studies the geometry
of the foliation of $\mathcal{U}$ by stable leaves. This is the
renormalization-theoretic underpinning of universality: two maps on the
same leaf are asymptotically indistinguishable under $\mathcal{R}^n$.

\subsection*{Status}
\conjectural{} The foliation theorem is known in the analytic category
(Lyubich~\cite{lyubich1999}); its extension to the specific family
$\{f_{\lambda,h}\}$ and the identification of which leaf this family
belongs to require explicit computation.

\depends{Section~5 of this volume: hyperbolicity;
Section~3 of this volume: $h \to 0$ limit}

%---------------------------------------------------------------
\section{Scale-Invariance of the TOGT Pipeline Under Renormalization}
%---------------------------------------------------------------

\subsection*{Description}
This section assembles the results of Part~II into a statement about the
TOGT pipeline: under iterated renormalization, the operator $\g_\lambda$
(via its discretization) converges to the universal fixed point $\g_\infty$,
and the large-scale structure of the pipeline
$S(L) \to \Phi \to \g \to \mathcal{K}(L) \to \text{Space}$
becomes independent of $\lambda$. This is the TOGT analogue of universality
in the classical sense. The extended pipeline
$S(L) \to \Phi \to \g_\lambda \xrightarrow{\mathcal{R}^n} \g_\infty \to
\mathcal{K}(L) \to \text{Space}$
proposed in Volume One (\S3.3) is given its first rigorous content here.

\subsection*{Status}
\conjectural{} Full rigour requires resolving the $h \to 0$ problem of
Section~3. The result is provable conditional on that resolution.

\depends{The full extended pipeline (V1, \S3.3);
Sections~4--6 of this volume: fixed point, hyperbolicity, stable manifolds}

% ============================================================
\part{The Topological Invariant 33}
% ============================================================

%---------------------------------------------------------------
\section{Candidate Definitions of the Topological Invariant}
%---------------------------------------------------------------

\subsection*{Description}
The scalar $m(\mathcal{M}(g^6(X))) = 33$ is stated as a conjecture in
Volume One (\S13) but the readout $m$ is not yet precisely defined.
This section surveys three candidate definitions: (i)~the sum of Betti
numbers of the sublevel-set filtration of the iterated state; (ii)~the
total persistence (sum of lifetimes in the persistence diagram); and
(iii)~a Morse-theoretic count of critical points of a smooth functional
associated to $g^6(X)$. Each candidate is evaluated against the
reproducibility protocol of Volume One (\S13.2) and the robustness data
(Table~1, V1).

\subsection*{Status}
\provable{} The definitions themselves require no new mathematics; the
contribution is selecting the one that matches the experimental data and
verifying that the selection is well-posed.

\depends{The Betti-sum protocol and robustness table (V1, \S\S13.2--13.3);
the honest-status remark (V1, \S13): $m$ requires a precise definition
before the conjecture is fully stated}

%---------------------------------------------------------------
\section{The Morse Complex of Iterated States}
%---------------------------------------------------------------

\subsection*{Description}
Given a chosen definition of $m$, this section constructs the Morse complex
$\mathcal{M}(g^k(X))$ for the iterated states produced by the TOGT
pipeline. It studies how the complex evolves under iteration --- how many
critical points appear or merge at each step --- and whether there is a
stabilisation phenomenon consistent with $m = 33$ at $k = 33$. The TTK
pipeline described in Volume One (\S13.2) is the computational model;
this section provides the mathematical framework it implicitly uses.

\subsection*{Status}
\conjectural{} The stabilisation at $m = 33$ is the empirical claim of
Volume One; whether it holds rigorously for the specific pipeline operators
$(\Comp, \Konst, \Fold, \Unfold)$ and any admissible $X$ is open. Proving
it would require control over the topology of the image of $g^k$ for all
admissible $X$.

\depends{The four-operator pipeline (V1, \S2);
the TTK protocol (V1, \S13.2);
Section~8 of this volume: candidate definition of $m$}

%---------------------------------------------------------------
\section{The Algebraic--Topological Separation, Revisited}
%---------------------------------------------------------------

\subsection*{Description}
The General Separation Theorem (V1, Theorem~12.2) proves that $33$ cannot
arise as an algebraic trace $\operatorname{Tr}(M_n^6)$ for any stable
$n < 33$ representation. This section revisits that result in light of the
topological invariant defined in Section~8: it asks whether $m = 33$ is
genuinely topological in the sense of being outside every algebraic
invariant of the linearized pipeline, and whether the separation is sharp
(i.e., whether there exists an $n = 33$ representation achieving the
bound). This sharpens the bridge between the two tiers of the theory.

\subsection*{Status}
\provable{} (sharpness of the bound) The bound $|\operatorname{Tr}(M^6)| \le n$
is attained when all eigenvalues equal $1$ (V1, Corollary~11.1); the
$n = 33$ case is therefore already provable from Volume One. The genuine
topological content of $m = 33$ is conjectural pending Section~9.

\depends{Theorem~12.2 and Corollary~11.1 (V1):
algebraic trace bound and separation;
Section~9 of this volume: the Morse complex}

%---------------------------------------------------------------
\section{Behaviour of the Topological Invariant Under Renormalization}
%---------------------------------------------------------------

\subsection*{Description}
This section addresses open problem~(ii) from Volume One (\S3.3): whether
the topological invariant $m(\mathcal{M}(g^6(X))) = 33$ is preserved,
transformed, or destroyed under $\mathcal{R}^n$. The heuristic computation
of Volume One (\S3.4) identified the degree-doubling assumption as the
critical gap; this section formulates that assumption precisely in terms of
the chosen definition of $m$ (Section~8) and states what a proof would
require.

\subsection*{Status}
\conjectural{} No proof strategy is currently available. The section
records the precise mathematical content of the open problem and the
conditions under which invariance, covariance, or transformation could be
established.

\depends{Section~8 of this volume: definition of $m$;
Sections~4--7 of this volume: renormalization dynamics;
The heuristic argument (V1, \S3.4): identifies the degree-doubling gap}

% ============================================================
\part{The $g^{64}$ Crystalline Spine}
% ============================================================

%---------------------------------------------------------------
\section{The $n = 64$ Operator Algebra}
%---------------------------------------------------------------

\subsection*{Description}
The Kether Orthogon conjecture (V1, Conjecture~14.1) posits that $g^{64}$
achieves a terminal crystalline fixed point with topological readout $33$.
A prerequisite is the $64$-dimensional operator algebra: the matrix
$M_{64}$ satisfying $\rho(M_{64}) = 1$, one contractive block, and one
expansive direction. This section constructs $M_{64}$ as a natural
extension of the canonical matrix $M_c$ (V1, \S11.1), verifies that the
algebraic trace bound gives $|\operatorname{Tr}(M_{64}^{64})| \le 64$, and
identifies which eigenvalue configurations would make $33 < 64$ a
consistent reading across both the algebraic and topological layers.

\subsection*{Status}
\provable{} The construction of $M_{64}$ and the trace bound follow
directly from Theorem~12.1 (V1) with $n = 64$; no new mathematics is
required beyond specifying the block structure.

\depends{Theorem~12.1 (V1): algebraic trace bound for general $n$;
the canonical matrix $M_c$ (V1, \S11.1);
Corollary~11.1 (V1): attainability of the bound}

%---------------------------------------------------------------
\section{Is $g^{64}$ a Fixed Point of $\mathcal{R}$?}
%---------------------------------------------------------------

\subsection*{Description}
The Kether Orthogon conjecture requires $g^{64}$ to be an attractor or
fixed point of the renormalization flow. This section formulates that
question precisely: does the renormalization orbit $\mathcal{R}^n(f_{\lambda,h})$
converge to $g^{64}$ (or to a map whose $64$-fold iterate has the conjectured
topological structure), or does it converge to $g_\infty$ (the Feigenbaum
fixed point) independently of iteration count? The two conjectures are
shown to be compatible only under specific conditions on the relationship
between $g_\infty$ and $g^{64}$.

\subsection*{Status}
\conjectural{} The relationship between $g_\infty$ (Feigenbaum fixed point)
and $g^{64}$ (crystalline spine) is not currently understood. This section
records the compatibility conditions as a precise open problem.

\depends{Conjecture~14.1 (V1): Kether Orthogon;
Sections~4--7 of this volume: renormalization dynamics;
Section~12 of this volume: $n = 64$ algebra}

%---------------------------------------------------------------
\section{The Orthogonality Condition}
%---------------------------------------------------------------

\subsection*{Description}
The Kether Orthogon conjecture also posits an orthogonality condition
$\langle \mathcal{M}(g^{64}(X)),\, g^{64}(X) \rangle_{\mathrm{TTK}} = 0$
(V1, \S14). Before this can be studied, the TTK inner product
$\langle \cdot, \cdot \rangle_{\mathrm{TTK}}$ on persistence modules must
be defined. This section surveys candidate inner products on the space of
persistence diagrams (e.g., the Wasserstein-based pairing, the kernel
inner product of Reininghaus et al.), evaluates each against the
orthogonality condition, and selects a working definition.

\subsection*{Status}
\conjectural{} The TTK inner product is explicitly listed as undefined in
Volume One (\S14); this section is the first step toward resolving that
gap. No orthogonality claim is made until a definition is fixed.

\depends{Conjecture~14.1 (V1): orthogonality condition;
Section~9 of this volume: Morse complex construction;
Section~10 of this volume: topological invariant under renormalization}

%---------------------------------------------------------------
\section{The $g^{64}$ Spine as Terminal Attractor: Synthesis}
%---------------------------------------------------------------

\subsection*{Description}
This closing section assembles the results of Part~IV into a unified
statement --- conditional on the conjectures of Sections~13 and~14 --- about
the terminal behaviour of the TOGT pipeline. It states what a complete
proof of the Kether Orthogon conjecture would require, which components
are now within reach (the trace bound, the $n = 64$ algebra), and which
remain genuinely open (the orthogonality condition, the relationship to
$g_\infty$). It also identifies the Lean~4 formalisation targets that
would close the remaining \texttt{sorry}s in the Volume~One nucleus.

\subsection*{Status}
\conjectural{} This section is explicitly synthetic; no new results are
claimed. It serves as a research agenda for the parts of Volume Two not
yet resolved.

\depends{All proved results in Volume One;
Sections~12--14 of this volume;
The Lean nucleus (V1, \S7): identifies the two scoped \texttt{sorry}s
that formalisation of this section would discharge}

% ============================================================
\section*{Summary of Dependencies and Status}
\addcontentsline{toc}{section}{Summary of Dependencies and Status}
% ============================================================

The table below summarises each section of Volume Two, its status, and
its primary dependency in Volume One.

\medskip
\noindent
\begin{tabular}{clll}
\hline
\S & Title (abbreviated) & Status & Primary V1 dependency \\
\hline
1  & Unimodal map spaces          & \textcolor{blue!70!black}{Provable}     & Prop.~3.2 \\
2  & $\mathcal{R}$ on function space & \textcolor{blue!70!black}{Provable}  & Prop.~3.2, \S3.3 \\
3  & Flow vs.\ discretization     & \textcolor{purple}{Conjectural} & Prop.~3.2, \S3.1 \\
4  & Feigenbaum fixed point       & \textcolor{blue!70!black}{Provable}     & Prop.~3.2 \\
5  & Hyperbolicity                & \textcolor{blue!70!black}{Provable}     & \S\S5--6 \\
6  & Stable manifolds             & \textcolor{purple}{Conjectural} & \S\S5--6, V2~\S5 \\
7  & Scale-invariance             & \textcolor{purple}{Conjectural} & \S3.3, V2~\S\S4--6 \\
8  & Candidate definitions of $m$ & \textcolor{blue!70!black}{Provable}     & \S\S13.2--13.3 \\
9  & Morse complex                & \textcolor{purple}{Conjectural} & \S2, \S13.2 \\
10 & Separation revisited         & \textcolor{blue!70!black}{Provable}     & Thm.~12.2, Cor.~11.1 \\
11 & $m$ under renormalization    & \textcolor{purple}{Conjectural} & \S3.4, V2~\S\S4--7 \\
12 & $n=64$ operator algebra      & \textcolor{blue!70!black}{Provable}     & Thm.~12.1 \\
13 & $g^{64}$ as fixed point      & \textcolor{purple}{Conjectural} & Conj.~14.1 \\
14 & Orthogonality condition      & \textcolor{purple}{Conjectural} & Conj.~14.1 \\
15 & Terminal attractor synthesis & \textcolor{purple}{Conjectural} & All of V1 \\
\hline
\end{tabular}

\medskip
\noindent
Of the fifteen sections, six have clear proof strategies using results
already established in Volume One (\S\S1, 2, 4, 5, 8, 10, 12). The
remaining nine depend on new definitions, external formalisation, or
unresolved foundational questions. The provable sections should be
completed first; they form the scaffold on which the conjectural sections
can be attempted.

\begin{thebibliography}{9}

\bibitem{collet1980}
P.~Collet and J.-P.~Eckmann.
\emph{Iterated Maps on the Interval as Dynamical Systems}.
Birkh\"{a}user, Boston, 1980.

\bibitem{lyubich1999}
M.~Lyubich.
Feigenbaum--Coullet--Tresser universality and Milnor's hairiness conjecture.
\emph{Ann.\ of Math.}\ \textbf{149} (1999), 319--420.

\bibitem{avila2006}
A.~Avila and G.~Forni.
Weak mixing for interval exchange transformations and translation flows.
\emph{Ann.\ of Math.}\ \textbf{165} (2006), 637--664.

\bibitem{avila2011}
A.~Avila and M.~Lyubich.
The full renormalization horseshoe for unimodal maps of higher degree:
exponential contraction along hybrid classes.
\emph{Publ.\ Math.\ Inst.\ Hautes \'Etudes Sci.}\ \textbf{114} (2011),
171--223.

\end{thebibliography}

\end{document}
